\documentclass[10pt,aspectratio=169]{beamer}
\usepackage{bma01}

\title[Asíntotas horizontales y oblicuas]{Límites al infinito}
\subtitle{Asíntotas horizontales y asíntotas oblicuas\\[2pt]\small Semana 4 -- Sesión 2}

\begin{document}

\begin{frame}
  \titlepage
\end{frame}

\begin{frame}{Objetivos de la sesión}
  \begin{enumerate}
    \item Definir $\lim_{x\to\pm\infty}f(x)=L$ y calcular límites al infinito.
    \item Resolver indeterminaciones $\frac\infty\infty$ e $\infty-\infty$ (dividir por la mayor potencia, racionalizar).
    \item Hallar \resaltar{asíntotas horizontales} y \resaltar{asíntotas oblicuas}.
    \item Combinar todas las asíntotas para esbozar gráficas.
  \end{enumerate}
  \vspace{1mm}
  {\small\textbf{Competencia 2 del sílabo:} Evalúa propiedades de límites y continuidad y su utilidad en aplicaciones de ciencia e ingeniería.}
  \begin{block}{Plan (2 horas; incluye verificación previa de 5 min)}
    \begin{tabular}{ll}
      0:00 -- 0:20 & Límites al infinito: definición y límites básicos\\
      0:20 -- 0:55 & Racionales, raíces y exponenciales. Asíntotas horizontales\\
      0:55 -- 1:35 & Asíntotas oblicuas. Esbozo con todas las asíntotas\\
      1:35 -- 1:45 & Aplicaciones (cinética enzimática)\\
      1:45 -- 2:00 & Actividad colaborativa, resumen y cierre metacognitivo
    \end{tabular}
  \end{block}
\end{frame}

\begin{frame}{Aula invertida: preparación y verificación}
  \small
  \begin{previa}[antes de la sesión]
    \begin{enumerate}
      \item Lea límites al infinito y asíntotas horizontales y oblicuas.
      \item Calcule $\lim_{x\to\infty}\frac{3x^2+1}{x^2-5}$ dividiendo por la mayor potencia.
      \item Pregunta guía: ¿puede una gráfica cortar a su asíntota horizontal?
    \end{enumerate}
  \end{previa}
  \begin{block}{Verificación al inicio (5 min, en pares)}
    \begin{enumerate}
      \item ¿Cuál es la asíntota horizontal de $\frac{2x+1}{x-3}$?
      \item ¿$\lim_{x\to-\infty}\frac{\sqrt{x^2}}{x}$ vale 1 o $-1$?
    \end{enumerate}
    Los pares discuten y votan; el docente aclara solo lo que genera dudas y ajusta el ritmo de la sesión.
  \end{block}
\end{frame}

\section{Límites al infinito}

\begin{frame}{Límites al infinito}
  \begin{columns}
    \begin{column}{0.55\textwidth}
      \begin{definicion}
        $\displaystyle\lim_{x\to+\infty}f(x)=L$ si
        \[ \forall\varepsilon>0\ \exists N>0:\ x>N\Rightarrow|f(x)-L|<\varepsilon. \]
        Análogo para $x\to-\infty$ ($x<-N$).
      \end{definicion}
      \begin{teorema}[límites básicos]
        Para $r>0$ racional: \ $\displaystyle\lim_{x\to+\infty}\frac1{x^r}=0$; \ $\lim_{x\to-\infty}\frac1{x^r}=0$ si $x^r$ está definido.
        \\ $\lim_{x\to+\infty}e^{-x}=0$, \ $\lim_{x\to+\infty}\arctan x=\frac\pi2$.
      \end{teorema}
      \small Las reglas algebraicas de límites siguen valiendo.
    \end{column}
    \begin{column}{0.45\textwidth}
      \centering
      \begin{tikzpicture}
        \begin{axis}[bma, xmin=0, xmax=10, ymin=0, ymax=2.5, width=6.5cm, height=5.3cm, xtick={4}, xticklabels={$N$}, ytick={1.3,1.6,1.9}, yticklabels={$L-\varepsilon$,$L$,$L+\varepsilon$}, samples=200]
          \fill[uniVerde!20] (axis cs:0,1.3) rectangle (axis cs:10,1.9);
          \fill[uniAzul!15] (axis cs:4,0) rectangle (axis cs:10,2.5);
          \addplot[curva, domain=0.2:10] {1.6+sin(deg(3*x))/x};
          \draw[dashed] (axis cs:0,1.6) -- (axis cs:10,1.6);
        \end{axis}
      \end{tikzpicture}

      \small Para $x>N$, la gráfica queda en la franja: puede incluso \emph{cruzar} $y=L$.
    \end{column}
  \end{columns}
\end{frame}

\begin{frame}{Racionales: dividir por la mayor potencia}
  \begin{ejemplo}[1]
    $\displaystyle\lim_{x\to\infty}\frac{3x^2-x+5}{2x^2+7}=\lim_{x\to\infty}\frac{3-\frac1x+\frac5{x^2}}{2+\frac7{x^2}}=\frac32$.
  \end{ejemplo}
  \begin{teorema}[comparación de grados]
    Si $P(x)=a_nx^n+\cdots$ y $Q(x)=b_mx^m+\cdots$, entonces cuando $x\to\pm\infty$:
    \begin{itemize}
      \item $n<m$: \ $\frac PQ\to0$.
      \item $n=m$: \ $\frac PQ\to\frac{a_n}{b_m}$ (cociente de coeficientes principales).
      \item $n>m$: \ $\frac PQ\to\pm\infty$ (el signo según $\frac{a_n}{b_m}$ y la paridad de $n-m$).
    \end{itemize}
  \end{teorema}
  \begin{ejemplo}[2]
    $\displaystyle\lim_{x\to-\infty}\frac{x^3+1}{2x-5}=+\infty$ \ (se comporta como $\frac{x^2}{2}$). \qquad
    $\displaystyle\lim_{x\to\infty}\frac{4x+1}{x^2-3}=0$.
  \end{ejemplo}
\end{frame}

\begin{frame}{Raíces: ¡cuidado con $x\to-\infty$!}
  \begin{ejemplo}[3]
    $f(x)=\dfrac{2x}{\sqrt{x^2+1}}$. \ Como $\sqrt{x^2}=|x|$:
    \[ \frac{2x}{|x|\sqrt{1+\frac1{x^2}}}\ \to\ \begin{cases}2,&x\to+\infty\ (|x|=x)\\-2,&x\to-\infty\ (|x|=-x)\end{cases} \]
    Dos asíntotas horizontales distintas: $y=2$ y $y=-2$.
  \end{ejemplo}
  \begin{ejemplo}[4: $\infty-\infty$, racionalizar]
    \[ \lim_{x\to\infty}\left(\sqrt{x^2+x}-x\right)=\lim_{x\to\infty}\frac{x}{\sqrt{x^2+x}+x}=\lim_{x\to\infty}\frac{1}{\sqrt{1+\frac1x}+1}=\frac12. \]
  \end{ejemplo}
  \small \textbf{Ejemplo 5.} $\lim_{x\to\infty}\frac{\sen x}{x}=0$ (sándwich: $\left|\frac{\sen x}x\right|\le\frac1x$), aunque la gráfica cruza infinitas veces $y=0$.
\end{frame}

\section{Asíntotas horizontales}

\begin{frame}{Asíntotas horizontales}
  \begin{definicion}
    La recta $y=L$ es una \resaltar{asíntota horizontal} de $f$ si $\lim_{x\to+\infty}f(x)=L$ ó $\lim_{x\to-\infty}f(x)=L$.
  \end{definicion}
  \begin{columns}
    \begin{column}{0.33\textwidth}
      \centering
      \begin{tikzpicture}
        \begin{axis}[bma, xmin=-6, xmax=6, ymin=-2.6, ymax=2.6, width=5cm, height=4.5cm, xtick=\empty, ytick={-2,2}, samples=200, domain=-6:6]
          \addplot[curva] {2*x/sqrt(x^2+1)};
          \draw[dashed, uniGranate] (axis cs:-6,2) -- (axis cs:6,2);
          \draw[dashed, uniGranate] (axis cs:-6,-2) -- (axis cs:6,-2);
        \end{axis}
      \end{tikzpicture}\\
      \small $\frac{2x}{\sqrt{x^2+1}}$: dos asíntotas.
    \end{column}
    \begin{column}{0.33\textwidth}
      \centering
      \begin{tikzpicture}
        \begin{axis}[bma, xmin=-6, xmax=6, ymin=-2, ymax=2, width=5cm, height=4.5cm, xtick=\empty, ytick={-1.5708,1.5708}, yticklabels={$-\frac\pi2$,$\frac\pi2$}, samples=200, domain=-6:6]
          \addplot[uniVerde, very thick] {rad(atan(x))};
          \draw[dashed, uniGranate] (axis cs:-6,1.5708) -- (axis cs:6,1.5708);
          \draw[dashed, uniGranate] (axis cs:-6,-1.5708) -- (axis cs:6,-1.5708);
        \end{axis}
      \end{tikzpicture}\\
      \small $\arctan x$.
    \end{column}
    \begin{column}{0.33\textwidth}
      \centering
      \begin{tikzpicture}
        \begin{axis}[bma, xmin=-1, xmax=12, ymin=-0.3, ymax=0.6, width=5cm, height=4.5cm, xtick=\empty, ytick=\empty, samples=200, domain=0.5:12]
          \addplot[curva2] {sin(deg(x))/x};
          \draw[dashed, uniAzul] (axis cs:-1,0) -- (axis cs:12,0);
        \end{axis}
      \end{tikzpicture}\\
      \small $\frac{\sen x}{x}$ cruza su asíntota.
    \end{column}
  \end{columns}
  \begin{observacion}
    Una función tiene a lo sumo \textbf{dos} asíntotas horizontales (una por cada extremo). La gráfica \textbf{puede cortarlas}.
  \end{observacion}
\end{frame}

\begin{frame}{Exponenciales y aplicación}
  \begin{ejemplo}[6]
    $\displaystyle f(x)=\frac{3e^x+1}{e^x+2}$: \ $x\to+\infty$: dividir por $e^x$: $\to3$; \ $x\to-\infty$: $e^x\to0$, $\to\frac12$.
    \ Asíntotas: $y=3$ y $y=\frac12$.
  \end{ejemplo}
  \begin{columns}
    \begin{column}{0.55\textwidth}
      \begin{ejemplo}[7: cinética enzimática]
        La ecuación de Michaelis--Menten da la velocidad de una reacción catalizada:
        \[ v([S])=\frac{V_{\max}[S]}{K_M+[S]}\ \xrightarrow[{[S]}\to\infty]{}\ V_{\max}. \]
        La asíntota horizontal es la velocidad máxima (saturación de la enzima). En $[S]=K_M$, $v=\frac{V_{\max}}2$.
      \end{ejemplo}
    \end{column}
    \begin{column}{0.45\textwidth}
      \centering
      \begin{tikzpicture}
        \begin{axis}[bma, xlabel={$[S]$}, ylabel={$v$}, xmin=0, xmax=10, ymin=0, ymax=1.2, width=6.3cm, height=4.6cm, xtick={1}, xticklabels={$K_M$}, ytick={0.5,1}, yticklabels={$\frac{V_{\max}}2$,$V_{\max}$}, samples=100, domain=0:10]
          \addplot[curva] {x/(1+x)};
          \draw[dashed, uniGranate] (axis cs:0,1) -- (axis cs:10,1);
          \draw[dotted] (axis cs:1,0) -- (axis cs:1,0.5) -- (axis cs:0,0.5);
        \end{axis}
      \end{tikzpicture}
    \end{column}
  \end{columns}
\end{frame}

\section{Asíntotas oblicuas}

\begin{frame}{Asíntotas oblicuas}
  \begin{columns}
    \begin{column}{0.55\textwidth}
      \begin{definicion}
        La recta $y=mx+b$ ($m\neq0$) es \resaltar{asíntota oblicua} de $f$ en $+\infty$ si
        \[ \lim_{x\to+\infty}\big[f(x)-(mx+b)\big]=0. \]
        (Análogo en $-\infty$.)
      \end{definicion}
      \begin{teorema}[cálculo]
        \[ m=\lim_{x\to\pm\infty}\frac{f(x)}{x},\qquad b=\lim_{x\to\pm\infty}\big[f(x)-mx\big], \]
        siempre que ambos límites existan y sean finitos.
      \end{teorema}
      \small Si $m=0$ se obtiene una asíntota horizontal.
    \end{column}
    \begin{column}{0.45\textwidth}
      \centering
      \begin{tikzpicture}
        \begin{axis}[bma, xmin=-4, xmax=6, ymin=-6, ymax=9, width=6.6cm, height=6cm, xtick={1}, ytick={2}, samples=300, restrict y to domain=-12:12, unbounded coords=jump]
          \addplot[curva, domain=-4:6] {(x^2+1)/(x-1)};
          \addplot[curva2, dashed, domain=-4:6] {x+1};
          \draw[gray, dashed] (axis cs:1,-6) -- (axis cs:1,9);
          \node[uniGranate] at (axis cs:4.6,3.6) {\scriptsize $y=x+1$};
        \end{axis}
      \end{tikzpicture}

      \small $f(x)=\frac{x^2+1}{x-1}$.
    \end{column}
  \end{columns}
\end{frame}

\begin{frame}{Ejemplos de asíntotas oblicuas}
  \begin{ejemplo}[8: racional, división]
    Si $\operatorname{gr}P=\operatorname{gr}Q+1$, la división da la asíntota:
    \[ \frac{x^2+1}{x-1}=x+1+\frac{2}{x-1}\quad\Rightarrow\quad y=x+1\ \text{(en ambos extremos)}. \]
  \end{ejemplo}
  \begin{ejemplo}[9: con raíces]
    $f(x)=\sqrt{x^2+2x}$. \ $m=\lim_{x\to+\infty}\frac{\sqrt{x^2+2x}}{x}=1$;
    \[ b=\lim_{x\to+\infty}\left(\sqrt{x^2+2x}-x\right)=\lim\frac{2x}{\sqrt{x^2+2x}+x}=1. \]
    En $+\infty$: $y=x+1$. \ En $-\infty$: $m=-1$, $b=-1$: $y=-x-1$. \ (¡Asíntotas distintas en cada extremo!)
  \end{ejemplo}
  \begin{ejemplo}[10]
    $f(x)=x+\dfrac{\sen x}{x}$ tiene asíntota oblicua $y=x$ y la corta infinitas veces.
  \end{ejemplo}
\end{frame}

\begin{frame}{Esbozo con todas las asíntotas}
  \begin{columns}
    \begin{column}{0.5\textwidth}
      \begin{ejemplo}[11]
        $f(x)=\dfrac{x^3}{x^2-4}$.
        \begin{itemize}
          \item Impar (simétrica respecto del origen).
          \item Verticales: $x=\pm2$. \ En $x\to2^+$: $\frac{8}{0^+}=+\infty$; \ $x\to2^-$: $-\infty$.
          \item División: $\frac{x^3}{x^2-4}=x+\frac{4x}{x^2-4}$ \ $\Rightarrow$ oblicua $y=x$.
          \item Corta a la asíntota oblicua solo si $\frac{4x}{x^2-4}=0$: en $x=0$.
        \end{itemize}
      \end{ejemplo}
    \end{column}
    \begin{column}{0.5\textwidth}
      \centering
      \begin{tikzpicture}
        \begin{axis}[bma, xmin=-6, xmax=6, ymin=-9, ymax=9, width=7cm, height=6.5cm, xtick={-2,2}, ytick=\empty, samples=400, restrict y to domain=-15:15, unbounded coords=jump]
          \addplot[curva, domain=-6:6] {x^3/(x^2-4)};
          \addplot[curva2, dashed, domain=-6:6] {x};
          \draw[gray, dashed] (axis cs:2,-9) -- (axis cs:2,9);
          \draw[gray, dashed] (axis cs:-2,-9) -- (axis cs:-2,9);
        \end{axis}
      \end{tikzpicture}
    \end{column}
  \end{columns}
\end{frame}

\begin{frame}{Actividad colaborativa: Saturación enzimática}
  \small
  \begin{activa}[ABP con datos]
    La velocidad de una reacción mide $v=1.75,\,2.80,\,4.00,\,5.09,\,5.89$ para $S=1,2,4,8,16$. Con el modelo $v=\frac{V_{\max}S}{K+S}$: \textbf{(a)} interprete $V_{\max}$ como asíntota horizontal; \textbf{(b)} estime $V_{\max}$ y $K$ ($K$ es el $S$ que da $V_{\max}/2$); \textbf{(c)} verifique el ajuste con una gráfica de los datos y la curva.
  \end{activa}
  \begin{itemize}
    \item \textbf{Organización:} equipos de 3 a 4 con roles rotativos (coordinador, redactor, verificador, portavoz); 12--15 minutos.
    \item \textbf{Producto:} Estimaciones de $V_{\max}$ y $K$, gráfica y discusión de la calidad del ajuste.
    \item \textbf{Retroalimentación:} el docente recorre los equipos, hace preguntas y no da la solución hasta la puesta en común.
  \end{itemize}
\end{frame}

\section{Ejercicios}

\begin{frame}{Ejercicios propuestos}
  \small
  \begin{ejercicio}
    \begin{enumerate}
      \item Calcule: (a) $\displaystyle\lim_{x\to\infty}\frac{(2x-1)^3}{(x+2)(4x^2+1)}$, \ (b) $\displaystyle\lim_{x\to-\infty}\frac{\sqrt{9x^2+1}}{3-x}$,
        \ (c) $\displaystyle\lim_{x\to\infty}\left(\sqrt{x^2+3x}-\sqrt{x^2-x}\right)$, \ (d) $\displaystyle\lim_{x\to\infty}\left(\sqrt[3]{x^3+x^2}-x\right)$.
      \item Halle todas las asíntotas: (a) $f(x)=\dfrac{2x^2-x}{x^2-1}$, \ (b) $g(x)=\dfrac{x^2-3x+2}{x+1}$, \ (c) $h(x)=\dfrac{x}{\sqrt{x^2-4}}$,
        \ (d) $k(x)=x\,e^{1/x}$, \ (e) $p(x)=\sqrt{4x^2+x}-x$.
      \item Esboce $f(x)=\dfrac{x^2}{x-2}$ con sus asíntotas.
      \item ¿Puede una función tener asíntota oblicua en $+\infty$ y horizontal en $-\infty$? Dé un ejemplo.
      \item La población de una especie sigue $P(t)=\dfrac{5000}{1+49e^{-0.3t}}$. Halle $P(0)$ y la asíntota horizontal. Interprete.
      \item Pruebe con la definición que $\displaystyle\lim_{x\to\infty}\frac{3x+1}{x+2}=3$.
    \end{enumerate}
  \end{ejercicio}
\end{frame}

\begin{frame}{Metacognición y uso responsable de la IA}
  \small
  \begin{columns}[t]
    \begin{column}{0.5\textwidth}
      \begin{metacog}[cierre, 3 min]
        \begin{enumerate}
          \metacomun
          \item ¿Qué significa ``tender a una asíntota'' en términos físicos? Explíquelo con este caso.
        \end{enumerate}
      \end{metacog}
    \end{column}
    \begin{column}{0.5\textwidth}
      \begin{iabox}
        Pida a una IA ajustar $V_{\max}$ y $K$ a los datos. Compare con su estimación manual y comente si el ajuste es razonable.\\[2pt]
        \textit{Regla:} primero intente usted; luego use la herramienta para contrastar, nunca para sustituir su razonamiento. Anote el cambio en su bitácora.
      \end{iabox}
    \end{column}
  \end{columns}
\end{frame}

\begin{frame}{Resumen}
  \begin{itemize}
    \item $\lim_{x\to\infty}f=L$: para $x$ grande, $f(x)$ está cerca de $L$. \ $\frac1{x^r}\to0$.
    \item Racionales: comparar grados; \ raíces: $\sqrt{x^2}=|x|$; \ $\infty-\infty$: racionalizar.
    \item \resaltar{Horizontal} $y=L$: $\lim_{x\to\pm\infty}f=L$ (a lo sumo dos).
    \item \resaltar{Oblicua} $y=mx+b$: $m=\lim\frac{f}{x}$, $b=\lim(f-mx)$; en racionales, por división.
    \item Las gráficas pueden cortar las asíntotas horizontales y oblicuas (incluso infinitas veces).
  \end{itemize}
  \vspace{3mm}
  \textbf{Próxima sesión:} continuidad y tipos de discontinuidad.
\end{frame}


\begin{frame}{Evaluación y trabajo posterior}
  \small
  \begin{evalbox}[Práctica del miércoles de la semana 4: PC1]
    \textbf{Cobertura:} semanas 1, 2 y 3. \textbf{Naturaleza:} práctica calificada individual escrita.
  \end{evalbox}
  \begin{itemize}
    \item \criterios
    \item \declaracionia
    \item \textbf{Portafolio:} guarde el producto de la actividad, la bitácora de metacognición y los ejercicios propuestos de esta sesión; son evidencia de progreso y base de la retroalimentación.
  \end{itemize}
\end{frame}

\end{document}
